\section{按诊断、图像和推断层级重读实验}
\label{sec:measurement_chain}
本节补足前述结果的观测链。“测量”也经过仪器响应和拟合；这里将\textbf{仪器信号/直接约束、派生量、机制解释}分别说明，避免把反演图当成逐点直接探测。除原文报告的误差外，不自行补造精度；以下分辨率不是所有测量量共有的误差条。

\subsection{Yoo2013/F679：按四张图还原MRX离子的经历}
\textbf{装置和诊断。}Yoo等（2013，\texttt{Yoo2013}；2014，\texttt{F679}）采用4.5 mTorr氦放电、无外加导场、磁通芯间距42 cm；后者给$d_i\simeq8$ cm。磁探针取得$\B$，其空间导数给$\J$，磁通的时间导数给面外重联电场；Langmuir探针给$n_e,T_e$。浮动电势$\Phi_f$经
\begin{equation}
\Phi_p\simeq\Phi_f+[3.3+0.5\ln(m_i/m_p)]T_e
\end{equation}
换算等离子体电势（此处$T_e$数值用eV、电势用V），再由$-\nabla\Phi_p$得到面内场。\textbf{电子动量平衡不是这一电势图的唯一来源}：作者另由$\vct u_e=\vct u_i-\J/(en)$、磁场和电子压强估计电场，积分后与探针换算电势比较；Yoo2013图2和F679图5的点线对照正是这个检验。\citep{Yoo2013,F679}

离子谱线He II 468.6 nm的中心位移约束视线流速，宽度约束温度；必须拟合13个精细结构分量，否则温度高估15--25\%。实际展示的较细流速图以经光谱标定的Mach探针为主。光谱空间平均约3--4 cm、曝光5.6 $\mu$s；磁探针最细径向0.6 cm、轴向3 cm。F679从4200余次放电筛出约30\%，每点平均7--15次，误差取样本标准差与诊断不确定度中较大者。因此二维图是重复放电拼接，有可重复性筛选，不是同一次放电的瞬时三维快照。

\textbf{读图顺序。}先读Yoo2013图2：势阱横向束缚、向下游变深变宽；再读图3(a)：离子穿越分离线后转向出流，图3(b)显示面内电功集中在下游。最后读图4：下游谱线既偏移又展宽；沿$Z$的流速先升、约$Z>9$ cm趋饱和时$T_i$才明显上升（图示由约数eV到约15--20 eV，非精确拟合数值）。温度横向峰比出流峰窄，因而不能只用$u_i^2$解释温度图。这是加速与温升空间分离的直接约束。

\textbf{哪些解释仍有竞争。}Yoo2013明确指出，所测热点不在经典无磁化黏性剪切最强处，经典热传导又难维持该梯度，故不能把这份MRX温升简单贴上“经典黏性加热”。原文提出再磁化延长驻留、束与背景碰撞、波粒散射、势阱反弹和磁通芯边界等可能性；上游离子平均自由程6--12 cm、下游2--5 cm支持下游碰撞有效。F679图9--10再以有/无现实量级离子碰撞的VPIC对照：两者均可有较大二阶矩，有碰撞时分布更接近Maxwell。\textbf{这份分布对照来自模拟；实验光谱的一维近Maxwell形状并不等于直接测得完整三维分布和熵。}同装置前后工作有数据继承，不能当作独立装置重复。
\loc{Yoo2013 / F679}{前者PDF第2--4页图2--4；后者第3--5页诊断与电势、第8--9页图9、10；第4页误差选取标准}

\subsection{MRX电子：电功集中处、温度峰和热量去向分开读}
F679图11显示电子由上游5--7 eV到分离线内11--12 eV，温度峰略在电子层外；图12的$\J_e\cdot\E$则集中于X点邻域，原文估计约70\%的电子沉积在小蓝框内。沿场传导把能量带离局地，所以温度最大处并非电功最大处。式(8)给$340\pm30$ kW电功，式(10)经典电阻约$59\pm12$ kW。约17\%的不足说明经典模型未闭合所测预算，却不能把剩余83\%直接分配给某一种波或“Hall电子热源”。\citep{F679}

\texttt{F674}（Yamada等，2015）把这种局地信息组成有限体积功率账本：$31.5\leq R\leq43.5$ cm、$0\leq Z\leq15$ cm，按环向对称积分；准稳态仍保留储存项，压力近似各向同性。局地关键量约10--15\%误差；热流采用Braginskii估计，氧杂质辐射由谱线与浓度估计，电荷交换损失因未测中性密度而取上界。图7数值因此混有测量、导数/积分、输运模型及上界，绝不是四个独立热量计读数。用输入Poynting功率作分母与用净转换磁功作分母也不同，且出流Hall场带走的Poynting通量不可丢弃。\citep{F674}
\loc{F679 / F674}{前者PDF第9--11页图11、12及式(6)--(19)；后者正式PDF第7页式(2)--(6)与损失估计，第8页图7，第9页表I}

\subsection{Yoo2017：大电场不一定意味着该处离子电功大}
Yoo等（2017，\texttt{Yoo2017Heating}）控制密度不对称：典型高/低密度侧约$4/0.5\times10^{13}\,\mathrm{cm^{-3}}$、重联场145/168 G，无外加导场。约2200次放电筛出35\%，每点7--10次，典型标准差10--15\%；离子光谱曝光10 $\mu$s、空间平均4 cm。图2显示强势垒偏向低密度分离线，使高密度侧来的多数离子留在出流内；图11因此给出一个反直觉但清楚的现象：\textbf{低密度侧面内场虽可约600 V/m，离子流向与电场未有效对齐，该处离子净电功仍小。}高密度侧场约150 V/m，离子电功约15 W/cm$^3$，也低于对称参照。电功需要$\J_i$与$\E$的点积，不能按场强排名。\citep{Yoo2017Heating}

图4电子热点却向低密度侧分离线延伸。作者逐项估计极化、梯度和曲率漂移，在该观测位置不足以解释温升，提出$E_\parallel$加速与有限碰撞的解释；$E_\parallel\sim40\pm20$ V/m是从较大总场投影的小量，误差不能省略。约4 cm平均自由程内的能量增量约2 eV，与观测量级相容。这是位置限定的机制估计，不能推广为所有区域曲率/Fermi通道都无效。

图8--9中的\emph{总电子电功}主要还含抗磁电流与垂直电场的功。原文式(7)--(10)用压力平衡、焓输运解释它，不能将其全部等同于上段单粒子的平行增能。本报告采用完整内能方程区分压力功和通量，亦不把一般压力功一概宣布为不可逆或“完全不加热”。图5的2D/3D VPIC给近似温度增量、不同Maxwell化程度；图6实验波动与温升弱相关（$r\simeq0.29,p\simeq0.086$），支持保留LHDI的可能作用，但未确立唯一因果。图10两物种功率更接近，既与电子低密度侧附加沉积有关，也与离子跨势垒路径减少有关，不能只说“非对称性使电子更热”。
\loc{Yoo2017Heating}{PDF第4页诊断，第5页图2，第7--10页图4--6及机制估计，第12--14页图8、9及式(7)--(10)，第15页图10、11；图10不是PDF第12页}

\subsection{MAST：先建立两个不同热点，再观察能量交换}
Tanabe等（2015，\texttt{Tanabe2015MAST}）图1中P3线圈驱动两环沿竖直方向合并，径向喷流进入下游闭合磁面。图2是单次放电8帧Thomson序列，0.1 ms间隔，电子温度从约10到200 eV发生在不到1 ms内，并与Mirnov磁信号相联系。图3的二维图则由多次放电移动采样位置拼接：电子窄峰宽2--5 cm，离子热点在下游更宽；不能把两类时间/空间成像方法混成单次全场电影。\citep{Tanabe2015MAST}

电子诊断为300通道Ruby和130通道YAG Thomson散射，含光学模糊的空间分辨约15/10 mm；离子为32弦Doppler层析，主要CVI 529.05 nm，早期冷帧用CIII。离子图的温度是杂质谱线反演，向主离子的解释需依赖碰撞与诊断适用性。图3离子热点颜色标度约百eV量级，不能拿L223高驱动模拟的keV峰冒充这一实验的实测值。图4给出离子下游双峰与电子中心峰在约4--10 ms交换时间后趋向三峰，图5在近固定$B_{rec}\simeq0.08$ T下改变导场0.4--0.8 T，电子中心峰更尖、所测下游离子温度基本不变。

作者给出两种使电子峰变尖的可能解释：更强导场降低跨场热输运、延长高电场区驻留，或减少meandering幅度使电流层更尖。缺乏直接磁探针，本文不能在两者之间作唯一判决。$B_{rec}$又由合并后30 ms的EFIT估计，并非每个加热时刻局地测量。对离子，强导场虽可抑制某些黏性阻尼，但闭合磁面延长驻留，仍可能最终耗散出流。\textbf{实验约束的是峰位、宽度、时序和参数趋势；没有直接给出局地黏性熵产、完整各向异性张量或普适两物种热化比例。}
\loc{Tanabe2015MAST}{PDF第2页诊断与几何，第3页图2--4，第4页图5及竞争解释；第5页原文引用[40]、[41]连接Stanier和Browning模拟}


\subsection{HIT-SI：高频驱动下的热离子与冷电子是哪一层证据}
Kaptanoglu等（2020，\texttt{Kaptanoglu2020}）第III.A节引用HIT-SI高频68.5 kHz测量，$T_e\sim5$--10 eV、$T_i\sim20$--30 eV；并提及\emph{另一个装置HIT-SI3}在14.5 kHz的初步Thomson电子温度5--10 eV。不能将后者写成HIT-SI同一次放电的共址双物种测量。作者主要比较光谱、干涉、磁探针的整体趋势；原始诊断学位论文及HIT-SI3会议结果本轮未独立核读，精确温度数值仅是该模型论文转述的量级约束。\citep{Kaptanoglu2020}

模型图2的密度是合成干涉弦平均；温度与图3功率大多是\emph{包含注入器}的体积量，图4才给空间切面。随着驱动频率增加，模型离子温度明显增加、电子较弱变化：场反转更频繁，黏性与压缩使离子增热，电子欧姆热增加的同时壁面热流也增加。温度是源和损失平衡，不能把低$T_e$解释为电子从未获得能量。装置持续1--2 ms，物种交换约100--200 $\mu$s，仍可慢于注入周期，支持独立两温方程。作者明确因缺少真实注入电路而避免严格定量验证；勘误正文也未取得，所以这里不从上述转述数值提炼精确加热比。
\loc{Kaptanoglu2020}{35页接收稿第5页式(16)、(17)与边界，第6--8页第III.A节及图2，第9页图3、4；第15页关于缺电路而不作定量验证的说明}

\subsection{新增空间观测参照：同一标度并不保证系数正确}
Drake等（2009，\texttt{Drake2009Pickup}）图6用Wind的22次磁剪切角大于$120^\circ$的出流穿越检验$\Delta T_p\propto(\Delta v_L)^2$；相对于式(8)的预测，拟合斜率只有0.39，不能只引用标度“吻合”而隐去系数偏低。图5示例的质子/氦核温升约5.5/20 eV，而预测约14.7/58.6 eV。图7及表I的ACE重离子用SWICS 12 min分辨率的径向$T_{rr}$，不是完整标量温度；第一次短穿越无居中采样，后两次相对局地磁场方向也不同。\citep{Drake2009Pickup}

这些是空间观测而非受控实验。作者把远下游边界结构、各向异性散射及膨胀列为解释缺口；本报告据此要求流体闭合同时检验总增量、各向异性和采样位置。此文不提供电子分配，也不能填补强导场2009年ApJ论文尚未核读的缺口。
\loc{Drake2009Pickup}{PDF第8--10页图5--7、表I，第11--12页讨论；SWICS只测径向温度的限制见第9页}
